% SOURCE_AUTHORITY: sga5_fr_workpass.tex, lines 10683--10714 (LF-normalized unit).
% SOURCE_SHA256: 791F4EFFC5E02832D5D77ED03518C8156D6F07E4C8238B03545DB93D883FBB28
Para el enunciado del corolario siguiente, se conviene en poner, para todo $X$-esquema $h:T\to X$, toda parte cerrada $Z$ de $T$ y todo complejo $M$ de $A_X$-módulos,
\[
H^*(T,M)=H^*(T,h^*(M)),\qquad
H^*_Z(T,M)=H^*_Z(T,h^*(M)).
\]

\begin{statement}{Corolario 8.3}
Para todo complejo $M$ de $A_X$-módulos y todo $p\in\mathbb Z$, las aplicaciones
\[
H^p(X,M)\oplus H^{p-2}(Y',M\otimes\mu^{\otimes -1})
\to
H^p(X',M)\oplus H^{p-2d}(Y,M\otimes\mu^{\otimes -d})
\tag{8.3.1}
\]
y
\[
H_Y^p(X,M)\oplus H^{p-2}(Y',M\otimes\mu^{\otimes -1})
\to
H_{Y'}^p(X',M)\oplus H^{p-2d}(Y,M\otimes\mu^{\otimes -d})
\tag{8.3.2}
\]
definidas por la matriz
\[
\begin{pmatrix}
f^*&j_*\\
0&g_*
\end{pmatrix}
\]
son biyecciones.
\end{statement}

\emph{Demostración.} Denotemos provisionalmente por $u$ el morfismo (8.1.1). Teniendo en cuenta la fórmula de proyección, las flechas (8.3.1) y (8.3.2) son las inducidas en los objetos de cohomología de grado $p$ por los morfismos $R\Gamma(X,u\otimes\operatorname{id}_M)$ y $R\Gamma_Y(X,u\otimes\operatorname{id}_M)$, respectivamente, de donde se deduce la afirmación.
